\section{Problem setting}
A scalar advection-diffusion problem without reaction.
\begin{equation} \label{eq:ad}
		  u_t + \div (F(u) - D(u) \nabla u) = 0, \ \bfx\in\Re^d
\end{equation}
The variable $D(u)$ will be treated with  Kirchhoff transformation.
The diffusion coefficient $D$ is treated as a constant here.
This will not affect the following discussion on numerical scheme.
\begin{equation} \label{eq:sad}
		  u_t + \div F(u) - D\Laplace u = 0, \ \bfx\in\Re^d
\end{equation}
\subsection{Finite volume scheme}
Replace scalar $u$ with cell averaged value 
\begin{equation}\label{eq:cellaveu}
		  \bar u = \frac{1}{|E|} \int_E u(\bfx, t) d\bfx
\end{equation}
and Lax-Friedrichs numerical flux
\begin{equation}\label{eq:lax}
		  f(u) = \frac{1}{2} [F(u^+)  + F(u^-) \cdot \nu_E - \alpha (u^+ - u^-)]
\end{equation}
The finite volume scheme of advection-diffusion problem is:
\begin{equation}\label{eq:fv}
		  \bar u_t + \frac{1}{|E|} \int_{\partial E} (\hat f(\bar u) - \nabla \bar u \cdot \nu_E) d\sigma(\bfx) = 0
\end{equation}
In 2D quadrilateral logically rectangular mesh, 
we apply gaussian quadrature every edge as the approximation of integral of flux. 
I use Multi-level WENO reconstruction scheme to reconstruct the point value in 
advection numerical flux.
\begin{equation}\label{eq:LaxFriedrichs}
		  \hat f(u) = \frac{1}{2} [F(\mathscr{R}^+(\{\bar u\}_+)) + F(\mathscr{R}^-(\{\bar u\}_-))\cdot \nu_E - 
		  \alpha(\mathscr{R}^+(\{\bar u\}_+) - \mathscr{R}^-(\{\bar u\}_-))] 
\end{equation}
Let $\{\bar u\}_+ = \{\bar u_i \mid i \in \text{WENO stencil} \quad I_+ \}$. 
Diffusion flux is different from its advection counterpart.
I reconstruct derivative across the edge by sampling along the normal direction.

\section{Computation of Jacobian}
Implicit method contains solving a non linear system and a linear subproblem.
I write the general non linear system as:
\begin{equation} \label{eq:nonlinear}
		  \bar u_t + \mathscr{F}(\bar u) - \mathscr{G}(\bar u) = 0
\end{equation}
where $\mathscr{F}(\bar u)$ represents advection part, and $\mathscr{G} (\bar u)$ represents diffusion part.
The semi discretized form is 
\begin{equation}
		  \bar u^{n+1} - \bar u^{n} + \Delta t [\mathscr{F} (\bar u^{n+1}) - \mathscr{G}(\bar u^{n+1})] = 0
\end{equation}
There are many well-researched numerical methods that can be used to solve such non linear system.
%I use non linear solver provided by PETSc.
%I provide routine computing right hand side and jacobian matrix.
%PETSc provides line search and preconditioner algorithm.
\subsection{ML-WENO reconstruction and derivatives}
I use ML-WENO reconstruction scheme to reconstruct point values
from cell-averaged value...
I use two stage non linear weighting. The first stage is
calculated as
\begin{equation}\label{eq:fstage}
		  \check \omega_l = \frac{\omega_l}{(\sigma_l + \epsilon h^2_0)^{\eta_l}} 
		  \quad \eta_l = \left \{\begin{aligned}
					 &1, \quad \text{if } r=1\\
					 &3, \quad \text{if } r=2\\
					 &4, \quad \text{if } r\geq 3
		  \end{aligned}
		  \right .
\end{equation}
\begin{equation}\label{eq:fstagenorm}
		  \bar \omega_l = \frac{\check \omega_l}{\sum_k \check \omega_k}
\end{equation}
The the second stage is calculated based on first stage:
\begin{equation}\label{eq:sstage}
		  \hat \omega_l = \frac{\bar \omega_l}{(\sigma_l + \epsilon h^2_0)^{sr_l}}
\end{equation}
\begin{equation}\label{eq:sstagenorm}
		  \tilde \omega_l = \frac{\hat \omega_l}{\sum_k \hat \omega_k}
\end{equation}
The reconstruction is expressed as the weighted linear combination of 
stencil polynomials.
\begin{equation}\label{eq:weno}
		  \mathscr{R}(\bfx) = \sum_{l\in L} P_l(\bfx) \tilde\omega_l(\bar u)
\end{equation}
Suppose the reconstruction happens in a given stencil $I$. 
The cell averaged values $\bar u$ in this stencil is represented as $\{\bar u_s\}_{s \in S}$.
The derivative of this reconstruction with respect to $\bar u_s$ is:
\begin{equation}\label{eq:derivR}
		  \frac{\partial \mathscr{R}}{\partial \bar u_s} = 
		  \sum_{l \in L} \Big(\frac{\partial P_l}{\partial \bar u_s}\tilde \omega_l + P_l\frac{\partial \tilde \omega_l}{\partial \bar u_s}\Big) 
\end{equation}
It consists of two parts, derivative of stencil polynomials $\frac{\partial P_l}{\partial \bar u_s}$ and 
derivative of non linear weights $\frac{\partial \tilde \omega_l}{\partial \bar u_s}$.
The stencil polynomial is a combination of single polynomials satisfying 
\begin{equation}\label{eq:spolyn}
		  \frac{1}{|E_k|} \int_{E_k} p_s(\bfx) d\bfx = \left \{ 
		  \begin{aligned}
					 &1 \quad s=k\\
					 &0 \quad s\neq k
		  \end{aligned}\right .
\end{equation}
I use tensorproduct polynomials for logically rectangular mesh.
The tensorporduct polynomial is written in the following form.
Coefficients $c_{\alpha}$ are determined by solving a linear
system satisfying (~\ref{eq:spolyn}).
\begin{equation}\label{eq:tensorProductPolyn}
		  p_s = \sum_{i < I} a^s_i (\frac{x-x_S}{h_S})^i \sum_{j < J} b^s_j (\frac{y-y_S}{h_S})^j = 
		  \sum_{\substack{\mid \alpha \mid \leq (I-1)\times (J-1)}} \tilde c^s_{\alpha} \Big(\frac{\bfx - \bfx_S}{h_S}\Big)^\alpha
\end{equation}
The stencil polynomial is expressed in terms of a 
combination of all these polynomials.
\begin{equation}\label{eq:Spolyn}
		  P_l(\bfx) = \sum_{s\in S} p^l_s(\bfx) \bar u_s
\end{equation}
Combining cell averaged values $\bar u_s$ with coefficients
\begin{equation}\label{eq:c_alpha}
		  c_{\alpha} = \sum_{s\in S} \tilde c^s_{\alpha} \bar u_s
\end{equation}
The derivative of polynomial is straightforward. The derivative of non linear weight is 
a little more complicated. Equation~\ref{eq:derivR} can be rewritten as:
\begin{equation}\label{eq:derivR_detail}
		  \frac{\partial \mathscr{R}}{\partial \bar u_s} = \sum_{l \in L} p^l_s(\bfx) \tilde \omega_l + 
		  \sum_{l\in L} P_l\frac{\partial \tilde \omega_l}{\partial \bar u_s}
\end{equation}
The derivatives of non linear weights are:
\begin{equation}\label{eq:derivNonlinearWgt}
		  \begin{aligned}
					 \frac{\partial \tilde \omega_l}{\partial \bar u_s} &= \frac{\hat \omega_{l,s}}{\sum_k \hat \omega_k} - 
		  \frac{\hat \omega_{l}\sum_k \hat \omega_{k,s}}{(\sum_k \hat \omega_{k})^2}\\
					 \frac{\partial \hat \omega_l}{\partial \bar u_s} &= \frac{\bar \omega_{l,s}}{(\sigma_l + \epsilon h_0^2)^{sr_l}} -
					 sr_l \frac{\bar \omega_l \sigma_{l,s}}{(\sigma_l + \epsilon h_0^2)^{sr_l+1}}\\
					 \frac{\partial \bar \omega_l}{\partial \bar u_s} &= \frac{\hat \omega_{l,s}}{\sum_k \hat \omega_k} - 
		  \frac{\check \omega_{l}\sum_k \check \omega_{k,s}}{(\sum_k \check \omega_{k})^2}\\
					 \frac{\partial \check \omega_l}{\partial \bar u_s} &= -\eta_l\sigma_{l,s}
					 \frac{\omega_l}{(\sigma_l + \epsilon h_0^2)^{\eta_l+1}}
		  \end{aligned}
\end{equation}
We calculate the derivatives of non linear weights with chain rule. 
We also need the derivative of smoothness indicator.
The Jiang-Shu smoothness indicator is:
\begin{equation}\label{eq:JiangShuSmoothnessInd}
		  \sigma_{JS} = \sum_{0 < \mid \beta \mid < r}\frac{h_0^{2\mid \beta \mid}}{\mid E_0 \mid} \int_{E_0} 
		  \Big ( \sum_{\mid \alpha \mid < r} c_{\alpha} \mathscr{D}^\beta \Big( \frac{\bfx - \bfx_S}{h_S}  \Big)^\alpha \Big )^2
		  d\bfx
\end{equation}
where order $r = (I-1) \times(J-1)$, a constant determined by the size of stencil.
We present an approximation of the Jiang-Shu smoothness indicator (cite ....) by 
omitting all corss product terms.
\begin{equation}\label{eq:PolynSmoothnessInd}
		  \begin{aligned}
					 \tilde \sigma_P &= \sum_{0 < \mid \alpha \mid < r} \sum_{0<\beta \leq \alpha}  
		  \Big( \frac{h_0}{h_S} \Big)^{2\mid \beta \mid} c^2_{\alpha} \Big(\frac{\alpha !}{(\alpha - \beta)!} \Big)^2
		  \frac{1}{E_0}\int_{E_0}\Big( \frac{\bfx - \bfx_S}{h_S} \Big)^{2(\alpha - \beta)}d\bfx \\
					 &= \sum_{0 < \mid \alpha \mid < r} 
					 \sum_{0<\beta\leq \alpha} \mathscr{C} (\alpha, \beta) c^2_{\alpha}
		  \end{aligned}
\end{equation}
the derivative of $\tilde \sigma_P$ can be written as:
\begin{equation}\label{eq:derivative_c_alpha}
		 \begin{aligned}
					\frac{\partial \tilde \sigma_P}{\partial \bar u_s} &=   \sum_{0 < \mid \alpha \mid < r} 
					 \sum_{0<\beta\leq \alpha} \mathscr{C} (\alpha, \beta) 2c_{\alpha}\frac{\partial c_{\alpha}}{\partial \bar u_s}\\
					 \frac{\partial c_\alpha}{\partial \bar u_s} &= \tilde c^s_{\alpha}
		 \end{aligned}
\end{equation}

\subsection{Derivative of flux}
Implicit method requires calculation of derivatives of both advection and diffusion flux.

\subsubsection{Derivative of advection flux}
Differentiate equation~\ref{eq:LaxFriedrichs} is straight forward.
\begin{equation}\label{eq:diffLaxFriedrichs}
		  \begin{aligned}
					 \frac{\partial \hat f(u)}{\partial \bar u_s} &= 
		  \frac{1}{2}
		  \Big[ \Big(\frac{\partial \mathscr{R}^+}{\partial \bar u_s} + \frac{\partial \mathscr{R}^-}{\partial \bar u_s}\Big)
		  \frac{\partial F}{\partial \bar u_s}\cdot \nu_E - \alpha 
		  \Big( \frac{\partial \mathscr{R}^+}{\partial \bar u_s}
			- \frac{\partial \mathscr{R}^-}{\partial \bar u_s}\Big)  \Big]\\
					 &=
					 \frac{1}{2} \Big[ 
					 \frac{\partial \mathscr{R}^+}{\partial \bar u_s}\Big( \frac{\partial F}{\partial \bar u_s}\cdot \nu_E 
					 - \alpha \Big) + 
\frac{\partial \mathscr{R}^-}{\partial \bar u_s}\Big( \frac{\partial F}{\partial \bar u_s}\cdot \nu_E 
					 + \alpha \Big) 
					 \Big]
		  \end{aligned}
\end{equation}

\subsubsection{Derivative of diffusion flux}



\section{Implicit time stepping}

\subsection{Backward Euler}

\subsection{Implicit SSP time stepping}

\section{Parallelism}
Parallel solver is inevitable in large scale computation.
We talk about mesh partitioning strategy 
acorss different processors in the following sections.
\subsection{Logically rectangular mesh and stencils}
Cells in logically rectangular mesh maintains logical 
order. Logically rectangular mesh shows more flexibility 
in some computation scenerios than cartesian mesh.
Figure~\ref{fig:fullmesh} 
shows a full view of a logically rectangular (quadrilateral) mesh.
\begin{figure}[htb]
		  \parbox{\textwidth}{\begin{minipage}[t]{0.45\textwidth}
		  \input sections/parallel_weno/mesh.tex
		  \caption{Logically rectangular 
		  mesh displayed in full. The index of 
		  a stencil (colored cells) is the 
		  same as the index of the red cell.}
\label{fig:fullmesh}
		  \end{minipage}
		  \hfill
		  \begin{minipage}[t]{.45\textwidth}
\input sections/parallel_weno/stencil.tex
					 \caption{Two sample stencils used by cell $E_{i,j}$.
					 The index of these two stencils are $S_{i-2,j}$ and $S_{i,j-2}$.
					 Consistent with the index of bottom left cell.}
		  \label{fig:stencil}
		  \end{minipage}}
\end{figure}
Counting all stencils without repeating is convenient 
in the logically rectangular mesh setting. 
We use the index of bottom left cell as the index of 
the corresponding stencil. A $\text{M}\times \text{N}$ mesh will then contain
$(\text{M}-\text{I}+1)\times(\text{N}-\text{J}+1)$ stencils 
of size $\text{I}\times \text{J}$.

\subsection{Mesh partition}
I examine the following partition of mesh without lossing generality.
The dashed line indicates ghost cells distributed by
MPI communication. 
The size of ghost layer is determined by the largest stencil
size required for cells in the ghost layer.
Increasing the order of reconstruction (increasing the size
of corresponding stencil) will thus increase the communication 
cost. However, the distribution of mesh only occur once 
in the beginning of the computation. Mesh will not change during
the process of compuation.
\begin{figure}\label{fig:partitionmesh}
\input sections/parallel_weno/quadmesh.tex
		  \caption{Illustration of mesh partition.
		  Ghost cells are drawn with dashed lines.
		  The size of ghost layer is determined 
		  by the number of cells need in the 
		  time stepping computation.}
\end{figure}
Figure~\ref{fig:partitionmesh} shows a stencil
covers ghost layer.

\subsection{Sparsity pattern of Jacobian matrix}

\subsection{Optimal preconditioner}

\begin{itemize}
					 \item 1. Incomplete LU decomposition
					 \item 2. Additive schwarz
\end{itemize}

\subsection{Scalability}

%\end{document}
